% ======================================================================
% INSERT BEFORE \section{Open questions} in chapter_d4_spectral.tex
% (i.e., before \label{sec:d4_open})
% ======================================================================

\section{The eigenvalue flow with compact momentum}
\label{sec:eigenvalue_flow}

The factorisation of Sec.~\ref{sec:factorisation} holds exactly at $k_4 = 0$.  At $k_4 \neq 0$, Systems~A and~B couple through terms proportional to $ik_4$ (Sec.~\ref{sec:chiral_mixing}).  The coupling preserves the block-diagonal structure of the transfer matrix at $k_2 = k_3 = 0$: the 20-component system separates into four independent blocks at \textit{all} values of~$k_4$, not only at $k_4 = 0$.  The blocks are:
\begin{center}
\small
\begin{tabular}{lccl}
\toprule
\textbf{Block} & \textbf{Size} & \textbf{System~A sector} & \textbf{System~B field} \\
\midrule
CG & $6 \times 6$ & Anti-plane $(u_3, \phi_2)$ $\to$ $\alpha$ & $\phi_{34}$ \\
DF & $6 \times 6$ & In-plane $(u_2, \phi_3)$ $\to$ $\alpha$ & $\phi_{24}$ \\
AE & $6 \times 6$ & Compression $(u_1)$ & $(u_4, \phi_{14})$ \\
Bt & $2 \times 2$ & Torsion $(\phi_1)$ & none \\
\bottomrule
\end{tabular}
\end{center}
Blocks CG and DF each produce one electromagnetic tunnelling eigenvalue; they are degenerate at $k_4 = 0$ (the two transverse polarisations of the photon) but may split at $k_4 \neq 0$.  Block~AE contains the compression sector (gravity) coupled to the compact displacement and propagation-direction mixed rotation.  Block~Bt (torsion) acquires a $k_4^2$ mass correction but no new coupling.

The transfer matrix $\mathbf{T}(k_4)$ for each block is computed by integrating the linearised Cosserat equations through one lattice period of the Peierls--Nabarro potential at transverse momentum $k_2 = k_3 = 0$ and compact momentum~$k_4$\footnote{Companion script \texttt{d4/d4\_transfer\_matrix\_20x20.py}~\cite{CoxGitHub}.}.  At $k_4 = 0$, the calibration reproduces the electromagnetic eigenvalue $|\lambda_{\text{CG}}| = |\lambda_{\text{DF}}| = \alpha$ to machine precision.

\paragraph{The flow.}  As $k_4$ increases from zero to the first Kaluza--Klein level $k_4\ell = 2\pi/3$, three effects are observed:

\textit{First}, the electromagnetic tunnelling amplitude decreases monotonically, from $|\lambda| = \alpha = 0.00730$ at $k_4 = 0$ to $|\lambda| = 0.00461$ at $k_4\ell = 2\pi/3$ --- a reduction of $36.8\%$.  The reduction is a geometric suppression: the $k_4^2$ mass terms in the Cosserat equations (from the four-dimensional Laplacian $\nabla_4^2 = \nabla^2 - k_4^2$) effectively raise the Peierls--Nabarro barrier, reducing the tunnelling rate.  This is not the renormalisation-group running of~$\alpha$ with energy; it is the dependence of the tunnelling amplitude on the compact momentum, a distinct physical variable.

\textit{Second}, the two electromagnetic eigenvalues split: $|\lambda_{\text{DF}}| > |\lambda_{\text{CG}}|$ at all $k_4 > 0$.  The splitting $\Delta|\lambda| = |\lambda_{\text{DF}}| - |\lambda_{\text{CG}}|$ is quadratic in $k_4$ at small $k_4$ (as required by time-reversal symmetry, which forbids odd powers) and reaches $\Delta|\lambda|/\alpha = 1.17\%$ at the first KK level.  The relative splitting $\Delta|\lambda|/\bar{\lambda}(k_4)$ grows from $0.01\%$ to $1.84\%$, demonstrating that the effect intensifies as the compact direction becomes more dynamically active.

\textit{Third}, the compression block (AE) remains opaque at all~$k_4$: the barrier is too high for single-node compression tunnelling regardless of the compact momentum.  The Born cluster mechanism that produces Newton's constant is robust.


\section{The chirality splitting: self-dual versus anti-self-dual}
\label{sec:chirality_splitting}

The splitting between blocks CG and DF has a clean structural origin.  In sector~C (anti-plane), the curl coupling enters as $-\kappa_c\phi_2'$ in the displacement equation, and the $ik_4$ mixing to the compact rotation enters as $-i\kappa_c k_4\phi_{34}/\mu_{\text{tot}}$.  In sector~D (in-plane), the curl coupling is $+\kappa_c\phi_3'$ and the $ik_4$ mixing is $+i\kappa_c k_4\phi_{24}/\mu_{\text{tot}}$.  The signs are correlated: the curl and the compact mixing have the \textit{same relative sign} in both sectors, but the absolute signs are opposite.

When the compact rotation is integrated out perturbatively (a round trip through the $ik_4$ coupling and back), the product of the two coupling insertions picks up a sign from the complex phases:
\begin{equation}
\text{CG:}\quad (-i)(+i) = +1\,,\qquad
\text{DF:}\quad (+i)(+i) = -1\,.
\label{eq:sign_product}
\end{equation}
The CG sector acquires a positive second-order shift (reduced tunnelling) and the DF sector acquires a negative shift (enhanced tunnelling), producing the splitting $|\lambda_{\text{DF}}| > |\lambda_{\text{CG}}|$.  The effect is second-order in $k_4$ (two coupling insertions) and first-order in $\kappa_c^2/(\mu_{\text{tot}}\gamma)$, consistent with the numerical observation.

In the self-dual / anti-self-dual language of Sec.~\ref{sec:chiral_decomposition}: the in-plane polarisation (DF, where the curl aligns with $ik_4$) is the self-dual component $\phi^+$, and the anti-plane polarisation (CG, where the curl opposes $ik_4$) is the anti-self-dual component $\phi^-$.  The lattice preferentially transmits the self-dual component.  The stacking direction $A \to B \to C$ defines the sense of ``alignment,'' and the opposite stacking $C \to B \to A$ would reverse the identification.  The vacuum crystal is optically active: it has a definite handedness for electromagnetic tunnelling above the Kaluza--Klein scale.


\section{The self-energy eigenvector and the compact-rotation leakage}
\label{sec:self_energy_leakage}

The electromagnetic eigenvector of the transfer matrix at $k_4 = 0$ lives entirely in the $(u, \phi_{\text{spatial}})$ subspace, with no component along the compact rotation $\phi_{\text{mixed}}$.  At $k_4 \neq 0$, the eigenvector acquires a compact-rotation component $v_{\text{mix}}$ through the $ik_4$ coupling.  The magnitude $|v_{\text{mix}}|^2$ measures the fraction of the electromagnetic mode that has ``leaked'' into the compact-direction rotation.

The leakage grows as $k_4^2$: $|v_{\text{mix}}|^2 = 0$ at $k_4 = 0$, rising to $|v_{\text{mix}}|^2 \approx 0.006$ ($0.6\%$) at $k_4\ell = 2\pi/3$.  This leakage is a first-order effect in the coupling strength $\kappa_c k_4/\mu_{\text{tot}}$: the compact rotation admixes linearly with $k_4$, and its probability (the square) grows quadratically.

\paragraph{The phase.}  The phase of $v_{\text{mix}}$ relative to $v_u$ is exactly $\pi/2$: the compact-rotation component is purely imaginary when the displacement component is real.  This is a direct consequence of the $ik_4$ coupling:  the factor of~$i$ rotates the compact-rotation amplitude by $90^\circ$ in the complex plane relative to the displacement.  The self-energy matrix element $\Sigma_{13} = \alpha(k_4)\,v_u\,v_{\text{mix}}^*$ is therefore \textit{purely imaginary}.

This has a physical interpretation of considerable importance.  A real self-energy matrix element ($\Sigma_{12}$) corresponds to a conservative, time-reversal-even coupling --- it shifts energy levels but does not distinguish forward from backward evolution.  An imaginary self-energy matrix element ($\Sigma_{13}$) corresponds to a dissipative, time-reversal-\textit{odd} coupling --- it distinguishes the direction of propagation through the compact circle and therefore breaks~$T$.

The compact direction is Euclidean time.  The imaginary phase of $\Sigma_{13}$ is the statement that the coupling between displacement and temporal rotation has a definite sense --- it distinguishes ``forward in time'' ($A \to B \to C$) from ``backward in time'' ($C \to B \to A$).  In the language of particle physics, this is \textit{CP violation}: a coupling that changes sign under the combined operation of charge conjugation and parity.

The monograph's existing CP violation arises from $\theta_{\text{ch}}^2 = \alpha^4/(4\pi^2)$, a real quantity that enters at fourth order through the centrosymmetric selection rule.  The $D_4$ mechanism provides CP violation at \textit{first order} through $\text{Im}(\Sigma_{13}) \propto \alpha(k_4)\,k_4$, with the imaginary unit supplied geometrically by the compact direction.  The two mechanisms are independent: $\theta_{\text{ch}}$ is a constitutive property of the elastic energy density (present at all temperatures), while $\text{Im}(\Sigma_{13})$ is a propagation property of the compact direction (present only when the compact direction is thermally active).


\section{The nested chirality and the thermal window}
\label{sec:nested_chirality}

The $D_4$ lattice possesses two independent sources of chirality, operating at different scales and with different symmetry structures.

\paragraph{Constitutive chirality: $\theta_{\text{ch}} = \alpha^2/(2\pi)$.}  This is the helical pitch of the Cosserat solid, derived from the centrosymmetric selection rule on the self-energy matrix (Sec.~\ref{sec:theta_derivation}).  It is real, time-reversal-even, and present at all temperatures.  It governs the weak force coupling, neutrino masses, the baryon asymmetry, and cosmic birefringence.  It is identical in the three-dimensional and four-dimensional readings of the lattice.

\paragraph{Propagation chirality: $\text{Im}(\Sigma_{13}) \propto \alpha\,k_4$.}  This is the optical activity of the compact direction, arising from the sign correlation between the curl coupling and the $\mathbb{Z}_3$ stacking (Sec.~\ref{sec:chirality_splitting}).  It is imaginary, time-reversal-odd, and present only when the compact direction is thermally excited ($T > T_{\text{geom}} = 23$~MeV).  It is specific to the $D_4$ interpretation.

The two chiralities are nested within a thermal hierarchy.  The compact direction's Matsubara frequencies are $k_4 = 2\pi n T/(\hbar c)$; the lowest nonzero mode ($n = 1$) becomes thermally active at $T \gtrsim T_{\text{geom}} = \hbar c/(3\ell\,k_B) \approx 23$~MeV.  As the temperature rises, the effective propagation chirality $\theta_{\text{eff}} \propto |v_{\text{mix}}|^2 \propto k_4^2 \propto T^2$ grows quadratically and overtakes the constitutive chirality $\theta_{\text{ch}}$ at a crossover temperature
\begin{equation}
T_{\text{cross}} \approx \frac{\hbar c}{\ell\,k_B} \times \frac{k_{4,\text{cross}}}{2\pi} \approx 103~\text{MeV}\,,
\label{eq:T_cross}
\end{equation}
where $k_{4,\text{cross}}\,\ell \approx 1.47$ is determined from the transfer matrix computation (the value of $k_4\ell$ at which $\theta_{\text{eff}} = \theta_{\text{ch}}$).  Above $T_{\text{cross}}$, the propagation chirality dominates.

This dominance is bounded above by the deconfinement temperature $T_c = 156.1$~MeV.  Above $T_c$, the temporal stacking order melts (the Polyakov loop disorders, Sec.~\ref{sec:polyakov_Tc}), the $\mathbb{Z}_3$ symmetry is restored, and the propagation chirality vanishes: the geometric handedness of the stacking sequence is destroyed along with the stacking itself.  The propagation chirality therefore exists only within a narrow thermal window:
\begin{equation}
T_{\text{geom}} \;<\; T \;<\; T_c \qquad (23~\text{MeV} < T < 156~\text{MeV})\,,
\label{eq:thermal_window}
\end{equation}
with the crossover at $T_{\text{cross}} \approx 103$~MeV dividing the window into a constitutive-dominated lower half and a propagation-dominated upper half.  The hierarchy $T_{\text{geom}} < T_{\text{cross}} < T_c$ is determined entirely by the lattice geometry and the Cosserat coupling number~$N^2 = 1/\pi$.


\section{Observational consequences}
\label{sec:d4_observational}

The propagation chirality is invisible at zero temperature: all laboratory experiments, all astrophysical observations except those probing the first microseconds of the universe or the interiors of neutron star mergers, operate at $T \ll T_{\text{geom}}$ and see only the constitutive chirality~$\theta_{\text{ch}}$.  Three environments reach the thermal window.

\paragraph{Heavy-ion collisions.}  The quark--gluon plasma produced at RHIC and LHC reaches temperatures of $300$--$500$~MeV, well above $T_c$.  At these temperatures, the temporal stacking is melted and the propagation chirality is absent.  As the fireball cools through the deconfinement transition and hadronises, the $\mathbb{Z}_3$ stacking recrystallises and the propagation chirality switches on.  The hadronisation epoch ($T \approx 156 \to 120$~MeV) passes through the upper half of the thermal window, where the propagation chirality dominates over~$\theta_{\text{ch}}$.

The chiral magnetic effect (CME) --- a charge current driven along an applied magnetic field by a chirality imbalance~\cite{Kharzeev2008,Fukushima2008} --- has been searched for extensively at RHIC and LHC through charge-dependent azimuthal correlations.  In the standard picture, the chirality imbalance arises from random topological fluctuations (sphalerons) and should average to zero over many events.  In the $D_4$ picture, the chirality imbalance during hadronisation is geometric (determined by the stacking handedness $A \to B \to C$) and has a \textit{definite sign} that does not cancel on average.  The two pictures make different predictions for the event-averaged charge-dependent correlator $\gamma$: zero (standard) versus nonzero with a definite sign ($D_4$).

\paragraph{The early universe.}  Between the electroweak transition ($T \sim 160$~GeV) and the QCD transition ($T \approx 156$~MeV), the universe cooled through a regime where the temporal stacking was disordered (above $T_c$) and the propagation chirality was absent.  The propagation chirality switched on at $T_c$ as the vacuum crystallised, and was active until the universe cooled below $T_{\text{geom}} \approx 23$~MeV (at $t \sim 10$~ms after the Big Bang).  The CP-violating, time-reversal-odd character of $\text{Im}(\Sigma_{13})$ satisfies Sakharov's third condition for baryogenesis~\cite{Sakharov1967} during this epoch.

\paragraph{Neutron star mergers.}  Binary neutron star merger remnants reach temperatures of $50$--$100$~MeV in the first milliseconds after coalescence, entering the lower half of the thermal window.  The propagation chirality would modify neutrino transport (through the effective weak coupling), the gravitational wave emission (through the microrotation sector), and the conditions for $r$-process nucleosynthesis in neutrino-driven winds.  The gravitational-wave birefringence ratio $\beta_{\text{GW}}/\beta_{\text{EM}} = 2/\pi$ predicted by the monograph for $T = 0$ (Sec.~\ref{sec:birefringence}) could be modified at merger temperatures by the propagation chirality enhancement.

\paragraph{Caveat.}  The quantitative predictions in this section are based on the transfer matrix eigenvalue flow, which gives the \textit{single-period} tunnelling amplitude as a function of~$k_4$.  The physical observables (the CME correlator, the baryogenesis source term, the neutrino opacity) involve the \textit{multi-period} Dyson equation that resums the transfer matrix into a propagator.  Whether the chirality splitting per period accumulates constructively or partially cancels upon resummation is an open calculation.  The structural result --- that the $D_4$ lattice has a time-reversal-odd chirality in the thermal window, with a definite sign and a first-order magnitude --- is established by the transfer matrix computation and does not depend on the resummation.
